\documentclass[10pt,a4paper]{amsart}
\usepackage[margin=19mm]{geometry}
\usepackage{amsmath,amssymb,mathtools}
\usepackage[scr=rsfs,cal=euler]{mathalfa}
\usepackage{xfrac}
\usepackage[unicode,hidelinks,bookmarks=false]{hyperref}
\newcommand{\ie}{i.e.\ }
\newcommand{\cf}{cf.\ }
\newcommand{\resp}{respectively\ }
\newcommand{\Subsection}{\subsection}
\providecommand{\nobreakdash}{\nobreak\hbox{-}}
\setlength{\emergencystretch}{2em}
\multlinegap0pt
\usepackage{bm}
\usepackage{bbm}
\usepackage{dsfont}
\usepackage{xspace}
\usepackage{cancel}
\usepackage{mathtools}
\usepackage{cleveref}
\usepackage{amsthm}
\makeatletter
\@ifundefined{theorem}{\newtheorem{theorem}{Theorem}}{}
\@ifundefined{proposition}{\newtheorem{proposition}[theorem]{Proposition}}{}
\makeatother
\DeclareMathOperator{\trace}{trace}
\newcommand{\bQ}{{\mathbf Q}}
\newcommand{\bS}{{\mathbf S}}
\newcommand{\bpi}{{\boldsymbol{\pi}}}
\newcommand{\bw}{{\mathbf w}}
\newcommand{\ip}[2]{\left\langle#1,#2\right\rangle}
\title[A Land of Oblique Duality for Frames and Probabilistic Frame]{A Land of Oblique Duality for Frames and Probabilistic Frames}
\author{Dongwei Chen, Emily J. King, Clayton Shonkwiler}
\date{Source: arXiv:2601.08028v1 (CC BY 4.0)}
\begin{document}
\maketitle

\noindent\emph{Source and licence.} A Land of Oblique Duality for Frames and Probabilistic Frames. Version arXiv:2601.08028v1 (CC BY 4.0), distributed under the Creative Commons Attribution 4.0 International licence, as recorded in the arXiv licence metadata for this exact version and shown on the abstract page https://arxiv.org/abs/2601.08028v1. Full rights notice: licenses/arxiv-2601.08028-CC-BY-4.0.txt.\par\smallskip

\noindent\textit{Editorial scope.} Counted unit: the Theorem whose source label is \texttt{thm:mixedCoherence} in the article (source0.tex lines 622--640, source label \texttt{thm:mixedCoherence}), delivered together with the article's own complete original proof of it (source0.tex lines 641--690, one \texttt{proof} environment of 50 lines). No mathematics is written, repaired or re-derived: statement and proof are the article's own text line for line. Required context retained uncounted: lemma:oblique\_dual\_potential (lines 549--555); eq:welch (lines 613--616). Every definition, display and label of those lines is retained unchanged. Omitted from this package and not redistributed: 1--548, 556--612, 617--621, 691--1459. Nothing inside a retained excerpt is omitted or altered: every retained body is delivered exactly as the article's lines are written, so no omission receipt of any kind accompanies this package. Citations: the retained excerpts carry no citation command at all, so no bibliography entry of the article is transcribed and no reference is redistributed. Reference adaptation: none was needed and none was made; every reference inside the delivered unit resolves inside it, so no unresolved reference or citation is produced. Printed slips kept literal: no printed word, symbol, hypothesis or number of the counted statement or of its proof was altered. Layout and class: the article's own class is \texttt{amsart} with packages including fontenc, amsfonts, mathrsfs, mathtools, amssymb, hyperref, cleveref, amsmath, blkarray, booktabs, bigstrut, graphicx, dutchcal, algorithm; the wrapper is the lane's pinned \texttt{amsart} editorial wrapper at 10pt on A4, which restates the theorem-like environments the retained text needs and adds the article's own macro definitions that the retained text needs (\texttt{\textbackslash bQ}, \texttt{\textbackslash bS}, \texttt{\textbackslash bpi}, \texttt{\textbackslash bw}, \texttt{\textbackslash ip}, \texttt{\textbackslash trace}), with extra wrapper packages (bm, bbm, dsfont, xspace, cancel, mathtools, cleveref). The delivered numbering is the unit's own. Assets: no retained span contains a figure environment, a graphics-inclusion command, a PDF-page inclusion, a table or a TikZ picture; no image file is copied into this package, the package declares no asset dependency and no image file is redistributed with it. Licence: version arXiv:2601.08028v1 of this article is distributed under the Creative Commons Attribution 4.0 International licence, as recorded in the arXiv licence metadata for this exact version and shown on the abstract page https://arxiv.org/abs/2601.08028v1. A full rights notice is delivered at licenses/arxiv-2601.08028-CC-BY-4.0.txt. Characters: every retained excerpt is pure ASCII, so the delivered engine, XeLaTeX with its default Latin Modern text font, reports no missing character.
\begin{proposition}\label{lemma:oblique_dual_potential}
    Let $\{\mathbf{w}_i\}_{i=1}^N$ be a frame for $W$ with frame operator $ \mathbf{S} = \sum_{i=1}^N \mathbf{w}_i\mathbf{w}_i^*$ and $\{\mathbf{v}_i\}_{i=1}^N$ an oblique dual frame of $\{\mathbf{w}_i\}_{i=1}^N$ on $V$, where $\mathbb{C}^n = W \oplus V^\perp$. Then 
    \begin{equation*}
         \sum_{i=1}^N \sum_{j=1}^N |\langle \mathbf{w}_i, \mathbf{v}_j \rangle |^2 \geq d_W,
    \end{equation*}
and the equality holds if and only if $\{\mathbf{v}_j\}_{j=1}^N$ is the canonical oblique dual $\{\bpi_{VW^\perp} \mathbf{S}^{\dagger} \mathbf{w}_j\}_{j=1}^N$. 
\end{proposition}

\begin{equation}\label{eq:welch}
    \max_{i \neq j} |\langle \mathbf{f}_i, \mathbf{f}_j \rangle|^2 
    \geq \frac{N-n}{n(N-1)},
\end{equation}

\begin{theorem}\label{thm:mixedCoherence}
     Suppose $W$ and $V$ are subspaces of $\mathbb{C}^n$ such that $\mathbb{C}^n = W \oplus V^\perp$. Let $\{\mathbf{w}_i\}_{i=1}^N$ be a frame for $W$ with frame operator $\mathbf{S}$ and $\{\mathbf{v}_i\}_{i=1}^N$ an oblique dual frame of $\{\mathbf{w}_i\}_{i=1}^N$ on $V$ such that $\langle \mathbf{w}_i, \mathbf{v}_i \rangle = \langle \mathbf{w}_j, \mathbf{v}_j \rangle$ for all $i$ and $j$. Then
    \begin{equation*}
       \underset{i \neq j}{\max} \ |\langle \mathbf{w}_i, \mathbf{v}_j \rangle |^2 \geq \frac{d_W(N-d_W)}{N^2(N-1)}.
    \end{equation*}
Furthermore, the equality holds if and only if any of the following equivalent conditions is true:
\begin{itemize}
    \item[$(1)$] For any $i \neq j$, $|\langle \mathbf{w}_i, \mathbf{v}_j \rangle |$ is constant, and $\mathbf{v}_j = \bpi_{VW^\perp} \mathbf{S}^\dagger \mathbf{w}_j$, for each $j$. 
    \item[$(2)$] The mixed Gram matrix $\mathbf{G} = (\langle \mathbf{w}_i, \mathbf{v}_j \rangle)_{ij}$ between $\{\mathbf{w}_i\}_{i=1}^N$ and $\{\mathbf{v}_i\}_{i=1}^N$ is
    \begin{equation*}
        \mathbf{G} =\frac{d_W}{N} \left(\mathbf{Id}_{N \times N} + \sqrt{\frac{N-d_W}{d_W(N-1)}}\mathbf{Q}\right),
    \end{equation*}
    where $\mathbf{Id}_{N \times N}$ is the identity matrix of size $N \times N$ and $\mathbf{Q}$ is a generalized signature matrix (self-adjoint, a zero diagonal and unimodular entries off the diagonal). 
    \item[$(3)$] For each $i$, define $
\boldsymbol{\psi}_i :=  \sqrt{\tfrac{N}{d_W}} \, (\mathbf{S}^\dagger)^{\frac{1}{2}} \mathbf{w}_i$.
Then $\{\boldsymbol{\psi}_i\}_{i=1}^N$ is an $(N,d_W)$- equiangular tight frame for $W$ and $\mathbf{v}_j = \bpi_{VW^\perp} \mathbf{S}^\dagger \mathbf{w}_j$, for each $j$.
    
\end{itemize}
\end{theorem}

\begin{proof}
Since $\{\mathbf{v}_i\}_{i=1}^N$ is an oblique dual of $\{\mathbf{w}_i\}_{i=1}^N$ on $V$ and the $\langle \mathbf{w}_i, \mathbf{v}_i \rangle$ are constant, we have 
$\langle \mathbf{w}_i, \mathbf{v}_i \rangle = \frac{d_W}{N}$ for each $i$ since
\begin{equation*}
   \sum_{i=1}^N \langle \mathbf{w}_i, \mathbf{v}_i \rangle = \trace\left( \sum_{i=1}^N \mathbf{w}_i \mathbf{v}_i^*\right)=\trace(\bpi_{WV^\perp}) =d_W. 
\end{equation*}
Therefore, 
    \begin{equation*}
    \begin{split}
        \underset{i \neq j}{\text{max}} \ |\langle \mathbf{w}_i, \mathbf{v}_j \rangle |^2 &\geq \frac{1}{N(N-1)} \sum_{i \neq j} |\langle \mathbf{w}_i, \mathbf{v}_j \rangle |^2 \\
        &= \frac{1}{N(N-1)} \left( \sum_{i=1}^N \sum_{j=1}^N | \langle \mathbf{w}_i, \mathbf{v}_j \rangle |^2 - \sum_{i=1}^N |\langle \mathbf{w}_i, \mathbf{v}_i \rangle |^2 \right) \\
        & \geq \frac{1}{N(N-1)} \left( d_W - \frac{d_W^2}{N} \right) = \frac{d_W(N-d_W)}{N^2(N-1)},
    \end{split}
    \end{equation*}
where the last inequality follows from \cref{lemma:oblique_dual_potential} and $d_W \geq \frac{d_W^2}{N}$ (since $d_W \leq N$).
Hence, the equality holds if and only if $|\langle \mathbf{w}_i, \mathbf{v}_j \rangle |$ is constant for all $ i \neq j$ and (by the equality condition in \cref{lemma:oblique_dual_potential}) $\mathbf{v}_j = \bpi_{VW^\perp} \mathbf{S}^\dagger \mathbf{w}_j$ for each $j$, where $\mathbf{S} = \sum\limits_{i=1}^N \mathbf{w}_i\mathbf{w}_i^*$. Therefore, equality is equivalent to $(1)$.

In turn, $(1)$ implies $|\langle \mathbf{w}_i, \mathbf{v}_j \rangle| =  \frac{1}{N}\sqrt{\frac{d_W(N-d_W)}{N-1}}$ for any $i \neq j$, so the mixed Gram matrix $\mathbf{G}$ between $\{\mathbf{w}_i\}_{i=1}^N$ and $\{\mathbf{v}_i\}_{i=1}^N$ can be written as 
    \begin{equation*}
        \mathbf{G} =\frac{d_W}{N} \left(\mathbf{Id}_{N \times N} + \sqrt{\frac{N-d_W}{d_W(N-1)}}\mathbf{Q}\right),
    \end{equation*}
    where $\mathbf{Id}_{N \times N}$ is the identity matrix of size $N \times N$ and $\mathbf{Q}$ is a matrix with zero diagonal and unimodular entries off the diagonal. Further, $(1)$ yields that the $i,j$ entry of the mixed Gram is 
    \begin{align*}
    \lefteqn{\ip{\bw_i}{\bpi_{VW^\perp}\bS^\dagger \bw_j} = \ip{\bpi_{WV^\perp}\bw_i}{\bS^\dagger \bw_j} = \ip{\bw_i}{\bS^\dagger \bw_j}}\\
    &= \ip{(\bS^{1/2})^\dagger\bw_i}{(\bS^{1/2})^\dagger \bw_j} =\overline{\ip{(\bS^{1/2})^\dagger\bw_j}{(\bS^{1/2})^\dagger \bw_i}} =\overline{\ip{\bw_j}{\bpi_{VW^\perp}\bS^\dagger \bw_i}},
    \end{align*}
    so $\bQ$ is self-adjoint and $(1)$ implies  $(2)$. 
Conversely, if $\mathbf{G}$ is given as in $(2)$, then for any $i \neq j$, $ |\langle \mathbf{w}_i, \mathbf{v}_j \rangle| = \frac{d_W}{N}\sqrt{\frac{N-d_W}{d_W(N-1)}}$. Hence
\begin{equation*}
       \underset{i \neq j}{\text{max}} \ |\langle \mathbf{w}_i, \mathbf{v}_j \rangle |^2 = \frac{d_W(N-d_W)}{N^2(N-1)},
    \end{equation*}
    which implies $(1)$. 

Finally, let us show $(2)$ and $(3)$ are equivalent. Define matrices $\mathbf{V}=(\mathbf{v}_1, \cdots, \mathbf{v}_N)$, $\mathbf{W} = (\mathbf{w}_1, \cdots, \mathbf{w}_N)$, and $\boldsymbol{\Psi} = (\boldsymbol{\psi}_1, \cdots, \boldsymbol{\psi}_N)$, where $\boldsymbol{\Psi} = \sqrt{\tfrac{N}{d_W}} \, (\mathbf{S}^\dagger)^{\frac{1}{2}} \mathbf{W}$. Thus, $\boldsymbol{\Psi}^* \boldsymbol{\Psi} = \tfrac{N}{d_W}  \mathbf{W}^* \mathbf{S}^\dagger \mathbf{W}$. If $(2)$ holds, then the equality holds and thus $\mathbf{v}_j = \bpi_{VW^\perp} \mathbf{S}^\dagger \mathbf{w}_j$ for each $j$. Hence  $\mathbf{V} = \bpi_{V W^\perp} \mathbf{S}^{\dagger}\mathbf{W} $ and
\begin{equation*}
    \frac{d_W}{N} \left(\mathbf{I}_{N \times N} + \sqrt{\frac{N-d_W}{d_W(N-1)}}\mathbf{Q}\right) = \mathbf{G} = \mathbf{V}^*\mathbf{W}  = \mathbf{W}^* \mathbf{S}^{\dagger} (\bpi_{W V^\perp}\mathbf{W})  = \mathbf{W}^* \mathbf{S}^{\dagger}\mathbf{W}.
    \end{equation*}
Therefore, the Gram matrix of $\{\boldsymbol{\psi}_i\}_{i=1}^N $ is given by
\begin{equation}\label{eqn:psiETF}
    \boldsymbol{\Psi}^* \boldsymbol{\Psi} = \tfrac{N}{d_W}  \mathbf{W}^* \mathbf{S}^\dagger \mathbf{W} = \mathbf{Id}_{N \times N} + \sqrt{\frac{N-d_W}{d_W(N-1)}}\mathbf{Q}
\end{equation}
and the Welch bound saturation condition for \eqref{eq:welch} implies that $\{\boldsymbol{\psi}_i\}_{i=1}^N$ is an $(N,d_W)$-equiangular tight frame. 
Conversely, if $(3)$ holds, then $\{\boldsymbol{\psi}_i\}_{i=1}^N$ is an $(N,d_W)$-equiangular tight frame and its Gram matrix is given as \cref{eqn:psiETF}. Since $\mathbf{v}_j = \bpi_{VW^\perp} \mathbf{S}^\dagger \mathbf{w}_j$ for each $j$, then  $\mathbf{V} = \bpi_{V W^\perp} \mathbf{S}^{\dagger}\mathbf{W} $ and thus
\begin{equation*}
\begin{split}
     \mathbf{G} = \mathbf{V}^*\mathbf{W} & =  \mathbf{W}^*  \mathbf{S}^{\dagger} (\bpi_{W V^\perp}\mathbf{W}) = \mathbf{W}^* \mathbf{S}^{\dagger}\mathbf{W} \\
     & = \frac{d_W}{N} \boldsymbol{\Psi}^* \boldsymbol{\Psi} =  \frac{d_W}{N} \left(\mathbf{Id}_{N \times N} + \sqrt{\frac{N-d_W}{d_W(N-1)}}\mathbf{Q} \right).
\end{split}
\end{equation*}
\end{proof}

\end{document}
